% !TEX program = xelatex
\documentclass{misisarticle}
\usepackage{flushend}

\grnti{28.17.19}
\asjc{2611, 1712}
\oecd{01.02}

\title{Математическая постановка и архитектура автономной библиотеки моделирования электрических цепей}

\misisauthor{Е.~А.~Федюнин}{кафедра инженерной кибернетики НИТУ «МИСИС»}{Москва, Россия}{egorfedunin@mail.ru}
\misisauthor{А.~И.~Тихонов}{ИГЭУ имени В.~И.~Ленина}{Иваново, Россия}{}

\abstract{Рассматривается перенос вычислительного ядра библиотеки моделирования электрических цепей ECLib из MATLAB в автономную библиотеку на C++. Формализованы графовое описание схемы, построение системы уравнений, шаг интегрирования и восстановление токов и напряжений. Предложена архитектура, в которой топологические блоки сохраняются между шагами при неизменных соединениях, а номиналы и источники обновляются отдельно. Для расчёта обмоток показано, как связать матрицу взаимных индуктивностей и резистивные потери с цепной моделью. Сформулирован критерий сравнения будущей реализации с MATLAB-прототипом по абсолютному и относительному допускам.}

\keywords{моделирование электрических цепей, дифференциально-алгебраические системы, графовые алгоритмы, разреженные матрицы, C++, верификация численных расчётов}

\newcommand{\solve}{\operatorname{solve}}
\newcommand{\eclibstage}[2]{%
  \fbox{\parbox{0.90\columnwidth}{\centering
    \strut\textbf{#1}\\[-0.1em]\footnotesize #2\strut}}}

\begin{document}
\maketitle

\section{Введение}

Для моделирования электрических машин, трансформаторов и реакторов часто требуется решить одну и ту же схему замещения при разных номиналах элементов и режимах источников. Библиотека ECLib автоматически формирует уравнения по описанию ветвей и их соединений~\cite{eclib2016}. Исходный MATLAB-прототип удобен для исследования алгоритма. Автономная реализация на C++ нужна, чтобы встраивать вычислительное ядро в инженерные приложения и выполнять серии расчётов без зависимости от среды MATLAB.

Цель работы --- формализовать переход от описания схемы к рассчитанным токам и напряжениям и спроектировать архитектуру, сохраняющую этот порядок вычислений. Рассматриваются построение системы, повторное использование неизменяемых данных и проверка переноса. Реализация ядра и измерение производительности относятся к следующему этапу.

\section{Анализ работ}

В исходной работе по ECLib~\cite{eclib2016} описана библиотека для виртуальных лабораторий численного исследования электрических машин. При переносе эталоном для сравнения результатов служит MATLAB-прототип. Другой подход к автоматическому построению уравнений представлен в модифицированном методе узловых потенциалов~\cite{ho1975}. В нём узловые потенциалы дополняются токами отдельных ветвей, что позволяет включать в общую систему, в частности, источники напряжения. Этот метод полезен для сопоставления способов выбора неизвестных; в базовой версии C++ сохраняется группировка ветвей MATLAB-прототипа.

Для больших схем существенны разреженность матрицы и порядок обработки её ненулевых элементов. Решатель KLU~\cite{klu2010} показывает, как блочная структура и разреженное LU-разложение используются в задачах моделирования цепей. В проектируемой библиотеке это обосновывает разделение сборки уравнений и их решения. Выбор между плотным и разреженным представлениями должен зависеть от структуры конкретной матрицы: взаимная индуктивность, например, может давать плотный блок. Одна из прикладных задач такого расчёта --- автоматизированное проектирование трансформаторов с использованием цифрового двойника~\cite{stulov2023}.

\section{Постановка задачи}

\subsection{Топология и входные данные}

На вход подаются матрица соединений, типы ветвей, матрица номиналов, заданные функции источников и начальные значения динамических переменных. Результат расчёта --- потенциалы узлов, токи и напряжения ветвей в исходной нумерации; для схем замещения обмоток дополнительно могут вычисляться потери.

Схема представляется ориентированным графом. Для каждой связной компоненты выбирается опорный узел, строка которого исключается из матрицы соединений. Пусть $A$ --- приведённая матрица соединений, $C_{\mathrm g}$ --- матрица независимых контуров, $i$ и $u$ --- векторы токов и напряжений ветвей, $\varphi$ --- вектор узловых потенциалов. При согласованной ориентации ветвей законы Кирхгофа~\cite{neyman1981} и связь напряжений с потенциалами записываются как
\begin{equation}
  A i=0,\qquad C_{\mathrm g}u=0,\qquad u=A^{\mathsf T}\varphi.
  \label{eq:kirchhoff}
\end{equation}

В алгоритме ECLib выделяются $y$-ветви, для которых расчётной величиной служит напряжение, и $z$-ветви, для которых расчётной величиной служит ток. К первой группе относятся источники ЭДС и ёмкости, ко второй --- сопротивления, индуктивности и источники тока. После выбора остовного леса внутри групп выделяются ветви дерева и связи. Внутренние неизвестные составляют напряжения $y$-ветвей дерева и токи $z$-связей.

\subsection{Система уравнений}

Пусть $D_y$ и $C_z$ --- блоки, определяемые топологией и порядком ветвей, а $Y_y$ и $Z_z$ --- блоки номиналов. В частности, диагональные блоки матрицы коэффициентов имеют вид
\begin{equation}
\begin{aligned}
  G_{11}&=D_yY_yD_y^{\mathsf T},\\
  G_{22}&=C_zZ_zC_z^{\mathsf T}.
\end{aligned}
\label{eq:blocks}
\end{equation}
При неизменной топологии $D_y$ и $C_z$ вычисляются один раз. Изменение номиналов требует пересборки зависящих от них коэффициентов, но не повторного поиска компонент и построения остова.

После подстановки известных источников дифференциальная и алгебраическая части системы записываются в виде
\begin{equation}
\begin{cases}
  G_p\dot X_p=F_p(t,X_p,X_r,e,j),\\
  G_rX_r+G_{rp}X_p=F_r(t,e,j).
\end{cases}
\label{eq:dae}
\end{equation}
Здесь $X_p$ содержит напряжения ёмкостей и токи индуктивностей, $X_r$ --- неизвестные резистивных ветвей; $e$ и $j$ --- значения источников ЭДС и тока. Матрицы $G_p$ и $G_r$ относятся к соответствующим подсистемам, а $G_{rp}$ связывает их. Далее рассматриваются схемы и режимы, для которых эти системы допускают однозначное решение.

\section{Описание методов}

\subsection{Численный шаг и восстановление величин}

В базовом переносе сохраняется явный шаг Эйлера, использованный в MATLAB-прототипе. На шаге $k$ сначала определяется согласованное алгебраическое состояние, затем вычисляется производная и обновляется динамический вектор:
\begin{equation}
\begin{aligned}
  X_r^k&=\solve(G_r^k,F_r^k-G_{rp}^kX_p^k),\\
  \dot X_p^k&=\solve(G_p^k,F_p^k),\\
  X_p^{k+1}&=X_p^k+\Delta t\,\dot X_p^k .
\end{aligned}
\label{eq:euler}
\end{equation}
Правая часть $F_p^k$ вычисляется после подстановки $X_r^k$. Оператор $\solve$ обозначает решение линейной системы через численное разложение без явного обращения матрицы~\cite{golub2013}. При новом значении $X_p$ алгебраическая часть пересчитывается. Физические величины восстанавливаются из законов Кирхгофа и соотношений элементов; обратные перестановки возвращают их в исходный порядок. Сохранение метода и шага времени позволяет сначала проверить эквивалентность двух реализаций.

\subsection{Применение к обмоткам}

Для схемы замещения обмотка разбивается на сегменты. При $i\ne j$ взаимная индуктивность определяется по геометрии пары витков~\cite{winding2015}:
\begin{equation}
  L_{ij}=W_iW_j\,M(r_i,z_i;r_j,z_j),\qquad i\ne j,
  \label{eq:inductance}
\end{equation}
где $W_i$ и $W_j$ --- числа витков, $r_i,z_i$ и $r_j,z_j$ --- радиальные и осевые координаты, а $M$ --- взаимная индуктивность пары единичных витков. Полученная матрица входит в блок индуктивностей цепной модели. Для множества резистивных ветвей $S_R$ мгновенная мощность потерь и среднее за период $T$ составляют
\begin{equation}
\begin{aligned}
  P_R(t)&=\sum_{j\in S_R} R_j i_j^2(t),\\
  \overline{P}_R&=\frac{1}{T}\int_{t_0}^{t_0+T}P_R(t)\,\mathrm{d}t.
\end{aligned}
\label{eq:losses}
\end{equation}
Здесь $R_j$ и $i_j$ --- сопротивление и ток ветви $j$, $t_0$ --- начало интервала усреднения.

\subsection{Архитектура и повторное использование данных}

В MATLAB-прототипе функции \nolinkurl{AnalisNet.m} и \nolinkurl{GetPrewMatrix.m} анализируют граф и строят постоянные блоки; \nolinkurl{GetPostMatrix.m} собирает текущую матрицу, \nolinkurl{StepCalcNet.m} выполняет шаг, а \nolinkurl{GetUI.m} восстанавливает физические величины. В версии C++ эти этапы предлагается разделить на модули (рис.~\ref{fig:architecture}).

\begin{figure}[!htb]
  \centering
  \begin{tabular}{@{}c@{}}
    \eclibstage{Входное описание цепи}{Соединения, номиналы, источники и начальные условия}\\[-0.15em]
    $\downarrow$\\[-0.15em]
    \eclibstage{Подготовка топологии}{Компоненты, остов, перестановки и постоянные блоки}\\[-0.15em]
    $\downarrow$\\[-0.15em]
    \eclibstage{Сборка системы и численный шаг}{Текущая матрица, правая часть и обновление состояния}\\[-0.15em]
    $\downarrow$\\[-0.15em]
    \eclibstage{Восстановление результата}{Токи, напряжения и потенциалы в исходном порядке}
  \end{tabular}
  \caption{Структура вычислительного ядра ECLib на C++}
  \label{fig:architecture}
\end{figure}

Топологический объект хранит перестановки и блоки из~\eqref{eq:blocks}. Рабочие буферы, текущие параметры и состояние интегратора принадлежат отдельному расчёту. Это позволяет разделять одну неизменяемую топологию между независимыми вариантами. При изменении только правой части допустимо повторное использование уже вычисленного численного разложения; при изменении коэффициентов его необходимо пересчитать. Если сохраняется структура ненулевых элементов и это поддерживается решателем, повторно используется символический анализ~\cite{klu2010}.

\subsection{Проверка эквивалентности}

Для проверки планируется использовать простые цепи с известным аналитическим решением, пример \nolinkurl{Test.m} из ECLib и расчётную модель \mbox{\texttt{Obmotka.xlsm}}. Для обеих реализаций задаются одинаковые начальные условия, источники и сетка времени. Сначала сопоставляются топология, матричные блоки и один шаг интегрирования; затем --- временные ряды токов и напряжений и рассчитанные потери. Для каждого сопоставляемого массива $r$ используется условие
\begin{equation}
  \|r_{\mathrm{C++}}^k-r_{\mathrm{MAT}}^k\|_\infty
  \le \varepsilon_a+\varepsilon_r\|r_{\mathrm{MAT}}^k\|_\infty ,
  \label{eq:tolerance}
\end{equation}
где $\varepsilon_a$ и $\varepsilon_r$ --- абсолютный и относительный допуски, а $\|\cdot\|_\infty$ --- максимальная норма. Допуски выбираются отдельно для каждой группы величин с учётом их масштаба и единиц измерения. Невязки законов Кирхгофа из~\eqref{eq:kirchhoff} позволяют дополнительно проверить восстановление результата в исходной нумерации.

Примеры сопоставления методов по явно указанным критериям приведены в работах студентов МИСИС~\mbox{\cite{podgorny2024,lichko2024}}. Время работы сравнивается для тех тестов, на которых расчёт C++ удовлетворяет условию~\eqref{eq:tolerance}. После проверки корректности планируется раздельно измерить подготовку топологии, сборку, решение системы и восстановление величин для скалярной и оптимизированной реализаций C++.

\begin{thebibliography}{9}

\bibitem{eclib2016}
Тихонов~А.~И., Стулов~А.~В., Дрязгов~Д.~Е.
Разработка автономной библиотеки моделирования электрических цепей для создания виртуальных лабораторий численного исследования моделей электрических машин //
Вестник ИГЭУ. 2016. Вып.~1. С.~34--39.
\href{https://doi.org/10.17588/2072-2672.2016.1.034-039}{DOI: \nolinkurl{10.17588/2072-2672.2016.1.034-039}}.

\bibitem{ho1975}
Ho~C.-W., Ruehli~A.~E., Brennan~P.~A.
The Modified Nodal Approach to Network Analysis //
IEEE Transactions on Circuits and Systems. 1975. Vol.~22, no.~6. P.~504--509.
\href{https://doi.org/10.1109/TCS.1975.1084079}{DOI: \nolinkurl{10.1109/TCS.1975.1084079}}.

\bibitem{klu2010}
Davis~T.~A., Palamadai Natarajan~E.
Algorithm 907: KLU, a Direct Sparse Solver for Circuit Simulation Problems //
ACM Transactions on Mathematical Software. 2010. Vol.~37, no.~3. Art.~36.
\href{https://doi.org/10.1145/1824801.1824814}{DOI: \nolinkurl{10.1145/1824801.1824814}}.

\bibitem{stulov2023}
Stulov~A., Tikhonov~A., Karzhevin~A.
Computer Aided Engineering System of Amorphous Transformers based on Digital Twin Technology //
RusAutoCon. 2023. P.~1046--1050.
\href{https://doi.org/10.1109/RusAutoCon58002.2023.10272939}{DOI: \nolinkurl{10.1109/RusAutoCon58002.2023.10272939}}.

\bibitem{neyman1981}
Нейман~Л.~Р., Демирчан~К.~С.
Теоретические основы электротехники. Т.~1. 3-е изд.
Л.: Энергоатомиздат, 1981. 536~с.

\bibitem{golub2013}
Golub~G.~H., Van Loan~C.~F.
Matrix Computations. 4th ed. Baltimore: Johns Hopkins University Press, 2013.

\bibitem{winding2015}
Тихонов~А.~И., Иванов~А.~В., Пайков~И.~А., Стулов~А.~В.
Математическое моделирование электромагнитных процессов в фольговых обмотках трансформаторов и токоограничивающих реакторов.
Иваново: ИГЭУ имени В.~И.~Ленина, 2015. 80~с.

\bibitem{podgorny2024}
Подгорный~Д.~А., Селезенёв~И.~А.
\href{https://misis.ru/files/-/7914c2bff86d1a41832cf9bc14e3de9b/kik_2024_1.pdf\#page=106}{Генерация оптического потока с помощью машинного обучения} //
Искусственный интеллект в промышленных, коммерческих, медицинских и финансовых приложениях.
М.: НИТУ «МИСИС», 2024. Вып.~2. С.~105--110.

\bibitem{lichko2024}
Личко~Д.~А., Рамзайцев~Д.~А.
\href{https://misis.ru/files/-/7914c2bff86d1a41832cf9bc14e3de9b/kik_2024_1.pdf\#page=112}{Распознавание поз нескольких объектов на изображении с использованием real-time моделей} //
Искусственный интеллект в промышленных, коммерческих, медицинских и финансовых приложениях.
М.: НИТУ «МИСИС», 2024. Вып.~2. С.~111--115.

\end{thebibliography}
\end{document}
